% 11장 결론. 세 문장으로 끝낸다.

\section{결론}

\textbf{첫째, 토지거래허가제는 거래를 확실하게 멈췄다.} 2020년 6월 지정 직후
3개월 거래량이 반사실 대비 $-71 \sim -80\%$다. 법정동 27개 가운데 4개를 처치로
고르는 17,550가지 중 진짜 배정이 1위이고, 사전 추세 위배를 실제 관측값의
1.6배까지 열어 두어도 구간이 0을 넘지 않으며, 합리적인 설계 조합 192가지가 전부
음수다. 완전히 다른 네 자치구에서 열 달 뒤에 일어난 2021년 지정에서도 같은
크기였고 두 사건의 차이는 0과 구분되지 않았다($p=0.715$).

\textbf{둘째, 그러나 모두에게 같지 않았고 그 차이는 제도 설계의 세부에서 왔다.}
2020년 지정은 대지 18\,m$^2$ 초과만 허가 대상이어서 소형 아파트는 거래가 줄지
않았다 --- 줄지 않은 정도가 아니라 통제군 대비 두 배 이상으로 늘었다. 2022년 6월
그 기준이 6\,m$^2$로 좁아지자 \emph{그때} 소형이 3분의 1 줄었고, 기준을 바꾸지
않은 2021년 재지정에서는 움직이지 않았다. 이 메커니즘에서 세운 예측 ---
``면적 기준이 6\,m$^2$인 2025년 재지정에서는 소형도 함께 줄어야 한다'' --- 이
결과를 보기 전에 기록되었고 자료에서 확인됐다(소형 $-71.5\%$).
\textbf{규제를 거느냐가 아니라 어떻게 거느냐가 효과를 정한다.}

\textbf{셋째, 가격을 잡았다는 증거는 약하다.} 평행추세를 가정하면 $-4 \sim -7\%$
이나 그 가정을 조금만 풀면 0과 구분되지 않고(붕괴값 0.22) 합성대조군으로도
확인되지 않는다($p=0.511$). 제도의 명분은 가격 안정인데 실제로 작동한 것은 거래
동결이었다. \textbf{이 어긋남을 숫자로 보인 것이 이 연구의 결론이다.}

\vspace{1em}

말할 수 없는 것도 함께 남긴다. 효과가 얼마나 오래 갔는지, 부과와 해제가 대칭인지,
대출 규제의 몫이 얼마인지는 이 설계로 답할 수 없다. \textbf{걸어 봤으나 가릴 수
없었다는 것도 결과이며}, 걸어 보지 않았으면 그 사실조차 몰랐을 것이다.

이 보고서의 또 하나의 축은 \textbf{분석의 전제가 되는 시스템}이었다. 원자료
수집부터 그림과 표 생성까지 48단계 파이프라인 하나로 돌고, 문서의 숫자가 분석
결과와 어긋나지 않는지·아무도 만들지 않는 산출물이 남아 있지 않은지·의존 목록이
실제 \texttt{import}와 맞는지를 매 실행마다 점검한다. 새 폴더에 내려받아 새
가상환경에서 돌려 결과 파일 25개가 바이트 단위로 같은 것을 확인했다. 그 과정에서
의존 목록의 구멍 둘과 석 달 묵은 유령 숫자를 찾았다.

\textbf{미시 자료로 내려간 것 자체가 기여는 아니다.} 기여는 미시 자료가 아니면
보이지 않는 것을 보였다는 데 있고, 면적 기준 발견이 그 주장의 가장 좋은 증거다.
동 평균 안에서 대형의 $-95\%$와 소형의 $+126\%$는 하나의 숫자로 뭉개진다.
\textbf{제목의 두 층 --- 미시 패널 구축과 인과효과 추정 --- 이 만나는 지점이
여기다.}
